\documentclass[11pt,a4paper]{article}
\usepackage[margin=2.5cm]{geometry}
\usepackage{booktabs}
\usepackage{amsmath}
\usepackage{array}

\title{\textbf{L1 Regularisation Values for Inverse-Design Surrogate Optimisation}}
\author{}
\date{Extracted: 2026-06-20}

\begin{document}
\maketitle
\thispagestyle{empty}

\section*{Extracted \texttt{lambda\_l1} Values}

The table below lists the L1 regularisation strength (\(\lambda\)) extracted from
every JSON file in \texttt{runtime\_data/surrogate\_optimisation/}.

\vspace{0.5em}
\begin{center}
\begin{tabular}{llll}
\toprule
\textbf{Target state} & \textbf{Node number} & \textbf{File} & \(\boldsymbol{\lambda_{\mathrm{L1}}}\) \\
\midrule
GHZ & 4 & \texttt{GHZ/4n.json} & $1\times10^{-4}$ \\
GHZ & 6 & \texttt{GHZ/6n.json} & $1\times10^{-5}$ \\
GHZ & 8 & \texttt{GHZ/8n.json} & $1\times10^{-6}$ \\
\midrule
W   & 4 & \texttt{W/4n.json}   & $1\times10^{-4}$ \\
W   & 6 & \texttt{W/6n.json}   & $1\times10^{-5}$ \\
W   & 8 & \texttt{W/8n.json}   & $1\times10^{-6}$ \\
\midrule
LC  & 4 & \texttt{LC/4n.json}  & $1\times10^{-4}$ \\
LC  & 6 & \texttt{LC/6n.json}  & $1\times10^{-5}$ \\
LC  & 8 & \texttt{LC/8n.json}  & $1\times10^{-6}$ \\
\bottomrule
\end{tabular}
\end{center}

\vspace{0.5em}
No files were missing and no \texttt{lambda\_l1} key was absent from any file.

\section*{Pattern}

The regularisation strength is \textbf{identical across all three target states}
(GHZ, W, LC) and decreases by one order of magnitude with each two-node increase
in system size:

\begin{itemize}
  \item 4-node: $\lambda = 1\times10^{-4}$
  \item 6-node: $\lambda = 1\times10^{-5}$
  \item 8-node: $\lambda = 1\times10^{-6}$
\end{itemize}

\section*{Suggested Methods Sentence}

\begin{quote}
``The L1 regularisation strength was set to
$\lambda = 10^{-4}$, $10^{-5}$, and $10^{-6}$
for the 4-, 6-, and 8-node systems, respectively,
applied uniformly across all target states (GHZ, W, and LC).''
\end{quote}

\end{document}
